% SPDX-License-Identifier: Apache-2.0
\section{Surface Syntax, Biased Toward Rust}
\label{sec:syntax}

Where the two languages disagree and Rust has an answer, Rust wins. Rust is
familiar, it is popular enough that engineers arrive already reading it, and
its constructs hold up under the coding patterns people now use. The bias
also earns its keep in a way that was not obvious in advance: it removes six
defects present in the sources, listed in Section~\ref{sec:defects}.

Where Rust has no answer, because the construct is about hardware and not
about software, the better of the two sources wins.

Rust's naming conventions apply throughout. Types, traits, buses,
transactions and modules take \code{CamelCase}. Values, functions, ports and
stages take \code{snake\_case}. Constants take
\code{SCREAMING\_SNAKE\_CASE}.

\subsection{The decision that does the most work}

TxHDL declares state with \code{var} and combinational logic with
\code{wire}. LHDL argues, in its eighth thesis, that the designer should not
make that distinction at all. Rust settles it with one keyword that neither
source has.

\begin{lstlisting}[caption={A binding and a place. \code{mut} is the whole of the distinction that \code{wire} and \code{var} used to draw.},label={lst:mut}]
let sum = a + b;            // a binding. Always a + b. A wire.
let mut count: u8 = 0;      // a place. Holds its value. A register.
\end{lstlisting}

\code{mut} states whether the value changes, which is a fact about the
design. It never states whether a flip-flop appears, which is a fact about
the implementation. This satisfies LHDL's eighth thesis and TxHDL's third
principle, ``variables become registers automatically based on usage'', with
one keyword readers already know. Both \code{var} and \code{wire} are
dropped.

\subsection{Conflicts that Rust settles}

Table~\ref{tab:syntax} lists them. Two entries are worth drawing out.

\emph{Bit slicing.} LHDL writes a Rust range and TxHDL writes a Verilog
descending pair. The range wins, which is the one place LHDL's spelling is
the more Rust-like of the two.

\emph{Awaiting.} Rust settled on postfix, and postfix chains where prefix
does not. Three of TxHDL's four await forms take it without argument. The
fourth, awaiting a condition, has no Rust ancestor, because Rust has no
cycle to wait for. The proposal is that a boolean is itself awaitable and
completes on the first cycle where it holds, which keeps one rule instead of
two. It reads oddly the first time, which is the argument against it.

\begin{table*}[t]
\caption{Conflicts Rust settles, with the precedent for each.}
\label{tab:syntax}
\centering
\footnotesize
\begin{tabular}{@{}llll@{}}
\toprule
Concern & LHDL & TxHDL & Merged \\
\midrule
Assignment          & \code{x := e}          & \code{x = e}             & \code{x = e} \\
Declaration         & \code{var x: T;}       & \code{var} / \code{wire} & \code{let} / \code{let mut} \\
Match arm           & \code{case P: \{..\}}  & \code{P => \{..\}}       & \code{P => \{..\}} \\
Match default       & \code{default:}        & \code{\_ =>}             & \code{\_ =>} \\
Enum variant        & \code{case Data(u32);} & \code{Add = 0x00,}       & both, comma separated \\
Struct fields       & \code{,} and \code{;}  & \code{;}                 & \code{,} \\
Struct literal      & none                   & \code{\{ .addr = 1 \}}   & \code{Req \{ addr: 1 \}} \\
Function            & \code{func f() -> T}   & \code{func f() -> T}     & \code{fn f() -> T} \\
Multi-cycle call    & none                   & \code{proc}              & \code{async fn} \\
Ternary             & none                   & \code{c ? a : b}         & \code{if c \{a\} else \{b\}} \\
Cast                & \code{u16(x)}          & \code{zext(x)}           & \code{x as u16} \\
Visibility          & \code{pub}             & none                     & \code{pub} \\
Packages            & \code{package a.b;}    & none                     & \code{mod}, \code{use a::b} \\
Generics            & \code{<type T>}        & \code{<T, const N: u32>} & \code{<T, const N: u32>} \\
Bit slice           & \code{w[0..8]}         & \code{w[7:0]}            & \code{w[0..8]}, \code{w[12..=14]} \\
Array type          & \code{u32[16]}         & \code{u32[16]}           & \code{[u32; 16]} \\
Loop                & \code{for (i in 0..n)} & \code{for (i in 0..n)}   & \code{for i in 0..n} \\
Loop step           & none                   & \code{step 4}            & \code{(0..n).step\_by(4)} \\
Labels              & none                   & \code{outer:}            & \code{'outer:} \\
Contract            & \code{implements}      & structural               & \code{trait}, \code{impl T for M} \\
Role                & \code{modport Master}  & \code{.master}           & \code{Bus::Master} \\
Assertion           & none                   & \code{assert c : "m"}    & \code{assert!(c, "m")} \\
Parallel branches   & \code{fork}/\code{join}& none                     & \code{join!(f(), g())} \\
Await               & none                   & \code{await port.req}    & \code{port.req.await} \\
Send                & none                   & \code{port.req <- \{..\}}& \code{port.req.send(..).await} \\
Callback            & none                   & \code{then (r) \{..\}}   & \code{.then(|r| \{..\})} \\
Timeout             & none                   & \code{timeout N else}    & \code{.timeout(N).await} \\
Tag application     & \code{@Sync expr}       & none                     & \code{\#[tag(Sync)]} \\
\bottomrule
\end{tabular}
\end{table*}

\subsection{Where Rust has no answer}

These constructs describe hardware, so the better source wins.

\begin{itemize}
  \item \term{Contract for wires}: \code{bus} with \code{channel}, from
        TxHDL. LHDL's \code{interface} does two jobs at once, a wire bundle
        and an implementation socket, which is why its grammar and its
        examples disagree. Splitting the jobs fixes both.
  \item \term{Contract for behaviour}: \code{trait}, replacing LHDL's
        interface in socket position. Late binding needs a name for a role
        an implementation fills, and that is what a trait is.
  \item \term{Timeless behaviour}: \code{flow}, from LHDL's dataflow model.
  \item \term{Sequenced behaviour}: \code{seq}, from TxHDL's sequence.
  \item \term{Synchronisation}: \code{tag}, from LHDL, spelled as an
        attribute, which states in one place what LHDL's grammar and
        examples spell three different ways.
  \item \term{Clock and reset}: \code{domain} in the configuration facet. A
        clock is a physical fact, and the design facet should not name one.
  \item \term{Late binding and foreign code}: \code{config}, \code{bind},
        \code{blackbox} and \code{foreign}, from LHDL. TxHDL has no
        mechanism at all.
  \item \term{Contiguous memory}: a named type, replacing LHDL's comma
        index. A comma inside brackets is too quiet for a decision that
        changes which primitive the synthesiser infers.
\end{itemize}
